% !TEX encoding = UTF-8
% !TEX Root = 0_main.tex
\vspace{-0.1cm}
\section{Rectified Adaptive Learning Rate}
\vspace{-0.1cm}
\label{sec:fix}

In the previous section, Equation~\ref{eqn:analytic-sqrt-var} gives the analytic form of $\Var[\psi(.)]$, where $\rho$ is the degree of freedoms. 
Here, we first give an estimation of $\rho$ based on $t$ to conduct a quantified analysis for $\Var[\psi(g_1, \cdots, g_t)]$, then we describe the design of the learning rate rectification, and  compare it to the heuristic warmup strategies.  

% \subsection{Estimation of $\rho$ and Variance}
\vspace{-0.1cm}
\subsection{Estimation of $\rho$}
\vspace{-0.1cm}

The exponential moving average (EMA) can be interpreted as an approximation to the simple moving average (SMA) in real application~\citep{nau2014forecasting}, \ie, 
\begin{align}
\vspace{-0.2in}
p\left(\frac{(1 - \beta_2)\sum_{i = 1}^t \beta_2^{t-i}g_i^2}{1-\beta_2^t} \right) \approx p\left(\frac{\sum_{i = 1}^{f(t, \beta_2)} g_{t+1-i}^2}{f(t, \beta_2)}\right).
\label{eqn:sma_ema}
\vspace{-0.2in}
\end{align}
where $f(t, \beta_2)$ is the length of the SMA which allows the SMA to have the same ``center of mass'' with the EMA.
In other words, $f(t, \beta_2)$ satisfies:
\begin{align}
\vspace{-0.7in}
% {(1 - \beta_2)}/{(1-\beta_2^t)} \sum_{i = 1}^t i \cdot \beta_2^{t-i} = 1 /{f(t, \beta_2)}\cdot\sum_{i = 1}^{f(t, \beta_2)} (t+1-i).
\frac{(1 - \beta_2)\sum_{i = 1}^t \beta_2^{t-i} \cdot i }{1-\beta_2^t} = \frac{\sum_{i = 1}^{f(t, \beta_2)} (t+1-i)}{f(t, \beta_2)}.
\label{eqn:sma_eql_ema}
\vspace{-0.7in}
\end{align}
% By solving the above equation, we have: 
By solving Equation~\ref{eqn:sma_eql_ema}, we have:
$f(t, \beta_2) = \frac{2}{1 - \beta_2} - 1 - \frac{2 t \beta_2^t}{1 - \beta_2^t}$.
In the previous section, we assume: $\frac{1-\beta_2^t}{(1 - \beta_2)\sum_{i = 1}^t \beta_2^{t-i}g_i^2} \sim \sics(\rho, \frac{1}{\sigma^2})$. 
Here, since $g_i \sim \gN(0, \sigma^2)$, we have $\frac{\sum_{i = 1}^{f(t, \beta_2)} g_{t+1-i}^2}{f(t, \beta_2)} \sim  \sics(f(t, \beta_2), \frac{1}{\sigma^2})$. 
Thus, Equation~\ref{eqn:sma_ema} views $\sics(f(t, \beta_2), \frac{1}{\sigma^2})$ as an approximation to $\sics(\rho, \frac{1}{\sigma^2})$.
Therefore, we treat $f(t, \beta_2)$ as an estimation of $\rho$. 
For ease of notation, we mark $f(t, \beta_2)$ as $\rho_t$. Also, we refer $\frac{2}{1 - \beta_2} - 1$ as $\rho_\infty$ (maximum length of the approximated SMA), due to the inequality $f(t, \beta_2) \leq \lim_{t \to \infty} f(t, \beta_2) = \frac{2}{1 - \beta_2} - 1$.


\vspace{-0.2cm}
\begin{algorithm}[t]
\DontPrintSemicolon
\KwIn{$\{\alpha_t\}_{t = 1}^T$: step size, 
$\{\beta_1, \beta_2\}$: decay rate to calculate moving average and moving 2nd moment,
$\theta_0$: initial parameter, $f_t(\theta)$: stochastic objective function.}
\KwOut{$\theta_t$: resulting parameters}
$m_0, v_0 \gets 0, 0$ (Initialize moving 1st and 2nd moment)\;
$\rho_\infty \gets 2/(1 - \beta_2) - 1$ (Compute the maximum length of the approximated SMA)\;
\While{$t = \{1, \cdots, T\}$}{
    $g_t \gets \nabla_{\theta} f_t(\theta_{t - 1})$ (Calculate gradients w.r.t. stochastic objective at timestep t)\;
    $v_t \gets \beta_2 v_{t - 1} + (1 - \beta_2) g_t^2$ (Update exponential moving 2nd moment)\;
    $m_t \gets \beta_1 m_{t - 1} + (1 - \beta_1) g_t$ (Update exponential moving 1st moment)\;
    $\widehat{m_t} \gets m_t / (1 - \beta_1^t)$ (Compute bias-corrected moving average)\;
    $\rho_t \gets \rho_\infty - 2t\beta_2^t/(1 - \beta_2^t)$(Compute the length of the approximated SMA)\;
    \uIf{the variance is tractable, \ie, $\rho_t > 4$}{
        $l_t \gets \sqrt{(1 - \beta_2^{t}) / v_t}$ (Compute adaptive learning rate)\;
        $r_t \gets \sqrt{\frac{(\rho_t - 4)(\rho_t - 2)\rho_\infty}{(\rho_\infty - 4)(\rho_\infty - 2)\rho_t}}$ (Compute the variance rectification term)\;
        $\theta_t \gets \theta_{t-1} - \alpha_t r_t \widehat{m_t} l_t$ (Update parameters with adaptive momentum)\;
    }
    \Else{
        $\theta_t \gets \theta_{t-1} - \alpha_t \widehat{m_t}$ (Update parameters with un-adapted momentum)\;
    }
}
\Return{$\theta_T$}
\caption{Rectified Adam. All operations are element-wise. }
% \vspace{-0.3cm}
\label{algo:cadam}
\end{algorithm}

\subsection{Variance Estimation and Rectification}
\vspace{-0.1cm}
Based on previous estimations, we have $\Var[\psi(.)] = \tau^2 (\frac{\rho_t}{\rho_t-2} - \frac{\rho_t \,2^{2\rho_t - 5}}{\pi}\gB(\frac{\rho_t-1}{2}, \frac{\rho_t-1}{2})^2)$.
The value of this function in the early stage is significantly larger than the late stage (as analyzed later, it decays roughly at the speed of $O(\frac{1}{\rho_t
})$). 
For example, 
% as in Figure~\ref{fig:var_approx}, 
the variance at $\rho_t = 5$ is over $100$ times larger than the variance at $\rho_t = 500$.
% \XD{Introducing Figure 7 requires explanation. How about just say we will analyze it in section 5.3.}
Additionally, based on Theorem~\ref{theorem: variance_mono}, we know $\min_{\rho_t} \Var[\psi(.)] = \Var[\psi(.)]|_{\rho_t = \rho_\infty}$ and mark this minimal value as $C_{\mbox{var}}$.
In order to ensure that the adaptive learning rate ($\psi(.)$) has consistent variance, we rectify the variance at the $t$-th timestamp as below,
% (with $\rho_t$), \ie,
\begin{align*}
% \vspace{-0.1in}
\Var[r_t \,\psi(g_1, \cdots, g_t)] = C_{\mbox{var}}
\quad
\mbox{where}
\quad
r_t = \sqrt{{C_{\mbox{var}}}/{\Var[\psi(g_1, \cdots, g_t)]}}.
% \label{eqn:analytic-rectification}
% \vspace{-0.1in}
\end{align*}
% \= \sqrt{\frac{(\rho_t - 4)(\rho_t - 2)\rho_\infty}{(\rho_\infty - 4)(\rho_\infty - 2)\rho_t}}.
Although we have the analytic form of $\Var[\psi(.)]$ (\ie, Equation~\ref{eqn:analytic-sqrt-var}), it is not numerically stable. 
Therefore, we use the first-order approximation to calculate the rectification term.
Specifically, by approximating $\sqrt{\psi^2(.)}$ to the first order~\citep{wolter2007taylor}, 
% Following~\citet{wolter2007taylor}, we have:
\begin{align*}
% \vspace{-0.2in}
\sqrt{\psi^2(.)} \approx \sqrt{\E[\psi^2(.)]} + \frac{1}{2\sqrt{\E[\psi^2(.)]}}(\psi^2(.) - \E[\psi^2(.)])
\quad\mbox{and}\quad
\Var[\psi(.)] \approx \frac{\Var[\psi^2(.)]}{4\E[\psi^2(.)]}.
% \label{eqn:sqrt-var}
% \vspace{-0.05in}
\end{align*}
Since $\psi^2(.) \sim \sics(\rho_t, \frac{1}{\sigma^2})$, we have:
% \vspace{-0.05in}
\begin{align}
% \vspace{-0.3in}
\Var[\psi(.)] \approx {\rho_t}/[{2(\rho_t - 2) (\rho_t - 4) \sigma^2}].
\label{eqn:variance_appro}
% \vspace{-0.1in}
\end{align}
% $
% \Var[\psi(.)] \approx {\rho_t}/[{2(\rho_t - 2) (\rho_t - 4) \sigma^2}].
% $
In Section~\ref{subsec:approx}, we conduct simulation experiments to examine Equation~\ref{eqn:variance_appro} and find that it is a reliable approximation. 
Based on Equation~\ref{eqn:variance_appro}, we know that $\Var[\sqrt{\psi(.)}]$ decreases approximately at the speed of $O(\frac{1}{\rho_t})$.
With this approximation, we can calculate the rectification term as:
% \vspace{-0.05in}
\begin{align*}
% \vspace{-0.3in}
r_t = \sqrt{\frac{(\rho_t - 4)(\rho_t - 2)\rho_\infty}{(\rho_\infty - 4)(\rho_\infty - 2)\rho_t}}.
% \vspace{-0.3in}
\end{align*}
% \vspace{-0.2in}
Applying our rectification term to Adam, we come up with a new variant of Adam, Rectified Adam (RAdam), as summarized in Algorithm~\ref{algo:cadam}. 
Specifically, when the length of the approximated SMA is less or equal than 4, the variance of the adaptive learning rate is intractable and the adaptive learning rate is inactivated. 
Otherwise, we calculate the variance rectification term and update parameters with the adaptive learning rate.
It is worth mentioning that, if $\beta_2 \leq 0.6$, we have $\rho_\infty \leq 4$ and RAdam is degenerated to SGD with momentum. 

\begin{figure}[t]
\begin{minipage}{0.69\textwidth}
% \vspace{-0.5cm}
\includegraphics[width=\textwidth]{fig/one_billion.pdf}
\vspace{-0.7cm}
 \captionof{figure}{Language modeling (LSTMs) on the One Billion Word.}
    \label{fig:1bw}
% \vspace{-0.4cm}
\end{minipage}
\begin{minipage}{0.3\textwidth}
    \centering
    % \begin{tabular}{c|c|c}
    %      Method &  Cifar10 & ImageNet\\
    %      \hline
    %      SGD   & 91.51 & 69.86 \\
    %      Adam   & 90.54 & 66.54 \\
    %      RAdam   & 91.38 & 67.62 \\
    % \end{tabular}
    \captionof{table}{Image Classification}
    \vspace{-0.2cm}
    \begin{tabular}{c|c|c}
        %  Method &  Cifar10 & ImageNet\\
        %  \hline
        & Method & Acc.\\
        \hline
        & & \\[-7pt]
      \parbox[t]{2mm}{\multirow{3}{*}{\rotatebox[origin=c]{90}{\small \;CIFAR10}}}
      & SGD & 91.51 \\[2pt]
        & Adam   & 90.54 \\[2pt]
        & RAdam   & 91.38 \\[2pt]
        \hline
        & & \\[-7pt]
        \parbox[t]{2mm}{\multirow{3}{*}{\rotatebox[origin=c]{90}{\small \;ImageNet}}}
        & SGD & 69.86 \\[2pt]
        & Adam & 66.54 \\[2pt]
        & RAdam & 67.62 
    \end{tabular}
% \end{table}
\end{minipage}
\end{figure}
\vspace{-0.4cm}
\begin{figure}[t]
\centering
\vspace{-0.2cm}
\includegraphics[width=\textwidth]{fig/cifa_imagenet.pdf}
\vspace{-0.6cm}
\caption{Training of ResNet-18 on the ImageNet and ResNet-20 on the CIFAR10 dataset.}
\vspace{-0.2cm}
    \label{fig:cifa10}
\end{figure}

\subsection{In Comparison with Warmup and Other Stabilization Techniques}
Different from the analysis in this paper,  warmup is originally proposed to handle training with very large batches for SGD~\citep{goyal2017accurate,gotmare2018a,bernstein2018signsgd,Xiao2017DSCOVRRP}.
We notice that $r_t$ has a similar form to the heuristic linear warmup, which can be viewed as setting the rectification term as $\frac{min(t, T_w)}{T_w}$. 
It verifies our intuition that warmup works as a variance reduction technique. 
% Comparing these two strategies, 
RAdam deactivates the adaptive learning rate when its variance is divergent, thus avoiding undesired instability in the first few updates. 
Besides, our method does not require an additional hyperparameter (\ie, $T_w$) 
% to control the variance reduction 
and can automatically adapt to different moving average rules. 

% In this paper, we identify and fix an underlying issue of adaptive optimization methods instead of neural architectures.
% In this paper
Here, we identify and address an underlying issue of adaptive optimization methods independent of (neural) model architectures.
Thus, the proposed rectification term is orthogonal to other training stabilization techniques such as  gradient clipping~\citep{bengio2013advances}, smoothing the adaptive learning rate (\ie, increasing $\epsilon$, applying geometric mean filter~\citep{chen2018closing}, or adding range constraints~\citep{luo2019adaptive}), initialization~\citep{balduzzi2017shattered,zhang2019fixup} and normalization~\citep{ba2016layer,ioffe2015batch}. 
% Indeed, these techniques can be integrated with our proposed variance rectification. 
Indeed, these techniques can be combined with the proposed variance rectification method.
%Specifically, since warmup is originally proposed to handle \emph{gradient variance} for SGD~\citep{goyal2017accurate,gotmare2018a,bernstein2018signsgd,xiao2019dscovr}, RAdam can also be integrated with the warmup heuristic to handle some extreme cases such as training with very large batches.


% To better understand the descent speed of the variance, we further derive a concise and reliable approximation. 
% Specifically, approximating $\sqrt{\psi^2(.)}$ to the first order~\citep{wolter2007taylor}, we have: 
% \begin{equation*}
% \sqrt{\psi^2(.)} \approx \sqrt{\E[\psi^2(.)]} + \frac{1}{2\sqrt{\E[\psi^2(.)]}}(\psi^2(.) - \E[\psi^2(.)])
% \quad\mbox{and}\quad
% \Var[\psi(.)] \approx \frac{\Var[\psi^2(.)]}{4\E[\psi^2(.)]}
% % \label{eqn:sqrt-var}
% \end{equation*}
% Since $\psi^2(.) \sim \sics(\rho_t, \frac{1}{\sigma^2})$, we have:
% \begin{equation}
% \Var[\psi(.)] \approx \frac{\rho_t}{2(\rho_t - 2) (\rho_t - 4) \sigma^2}.
% \label{eqn:variance_appro}
% \end{equation}
% In Section~\ref{subsec:approx}, we conduct simulation experiments to examine Equation~\ref{eqn:variance_appro} and find it is a reliable approximation. 

% \section{Rectified Adaptive Learning Rate}


% In the previous section, Equation~\ref{eqn:analytic-sqrt-var} gives the analytic form of $\Var[\psi]$, where $\rho$ is the degree of freedoms. 
% Here, we first give an estimation of $\rho$ and a quantified analysis of $\Var[\psi(.)]$; then we proceed to the design of the learning rate rectification and discussions on warmup strategies.  


% \subsection{Variance Estimation}

% As the exponential moving average (EMA) is widely used in eccomincs, it is usually interpretted as an approximation to the simple moving average (SMA)~\citep{nau2014forecasting}, \ie, 
% \begin{equation}
% \frac{(1 - \beta_2)\sum_{i = 1}^t \beta_2^{t-i}g_i^2}{1-\beta_2^t} \approx \frac{\sum_{i = 1}^{f(t, \beta_2)} g_{t+1-i}^2}{f(t, \beta_2)}.
% \label{eqn:sma_ema}
% \end{equation}
% where $f(t, \beta_2)$ is the length of the SMA which allows the SMA has the same ``center of mass'' with the EMA.
% In other words, $f(t, \beta_2)$ satisfies:
% $$
% \frac{(1 - \beta_2)\sum_{i = 1}^t \beta_2^{t-i} (t+1-i)}{1-\beta_2^t} = \frac{\sum_{i = 1}^{f(t, \beta_2)} (t+1-i)}{f(t, \beta_2)}.
% $$
% % Simplifies the above equation, we have:
% % $$
% % \frac{1}{1-\beta_2} - \frac{t  \beta_2^t}{1 - \beta_2^t}= \frac{f(t, \beta_2)+1}{2}
% % \quad\mbox{and}\quad 
% % f(t, \beta_2) = \frac{2}{1 - \beta_2} - 1 - \frac{2 t \beta_2^t}{1 - \beta_2^t}.
% % $$
% From this equation, we have: $f(t, \beta_2) = \frac{2}{1 - \beta_2} - 1 - \frac{2 t \beta_2^t}{1 - \beta_2^t}$.
% Since $g_i \sim \gN(0, \sigma^2)$, $\frac{\sum_{i = 1}^{f(t, \beta_2)} g_{t+1-i}^2}{f(t, \beta_2)} \sim  \sics(f(t, \beta_2), \frac{1}{\sigma^2})$. 
% % Since $\frac{\sum_{i = 1}^{f(t, \beta_2)} g_{t+1-i}^2}{f(t, \beta_2)}$ accords to the scaled inverse chi-square distribution with $f(t, \beta_2)$ degrees of freedoms, 
% % $
% %     \frac{1-\beta_2^t}{(1 - \beta_2)\sum_{i = 1}^t \beta_2^{t-i}g_i^2} \sim \sics(\rho, \frac{1}{\sigma^2})
% %     $
% %     is an approximation of
% %     $
% %     \frac{f(t, \beta_2)}{\sum_{i = 1}^{f(t, \beta_2)} g_{t+1-i}^2} \sim \sics(f(t, \beta_2), \frac{1}{\sigma^2}).
% %     $
% Additionally, in the previous section, we assume $\frac{1-\beta_2^t}{(1 - \beta_2)\sum_{i = 1}^t \beta_2^{t-i}g_i^2} \sim \sics(\rho, \frac{1}{\sigma^2})$. 
% Thus, Equation~\ref{eqn:sma_ema} views a sample from $\sics(f(t, \beta_2), \frac{1}{\sigma^2})$ as the approximation to a sample from $\sics(\rho, \frac{1}{\sigma^2})$.
% Therefore, we treat $f(t, \beta_2)$ as an estimation of $\rho$. 
% For ease of notation, we mark $f(t, \beta_2)$ as $\rho_t$. Also, we record $\frac{2}{1 - \beta_2} - 1$ as $\rho_\infty$ (maximum length of the approximated SMA), due to the inequality $f(t, \beta_2) \leq \lim_{t \to \infty} f(t, \beta_2) = \frac{2}{1 - \beta_2} - 1$.
% Then, we have $\Var[\psi(.)] = \frac{\rho_t}{\rho_t-2} - \frac{\rho_t \,2^{2\rho_t - 5}}{\pi}\gB(\frac{\rho_t-1}{2}, \frac{\rho_t-1}{2})^2$.
% To better understand the speed of the descent speed of variance, we further derive a concise and reliable approximation. 
% Specifically, approximating $\sqrt{\psi^2(.)}$ to the first order~\citep{wolter2007taylor}, we have: 
% \begin{equation*}
% \sqrt{\psi^2(.)} \approx \sqrt{\E[\psi^2(.)]} + \frac{1}{2\sqrt{\E[\psi^2(.)]}}(\psi^2(.) - \E[\psi^2(.)])
% \quad\mbox{and}\quad
% \Var[\psi(.)] \approx \frac{\Var[\psi^2(.)]}{4\E[\psi^2(.)]}
% % \label{eqn:sqrt-var}
% \end{equation*}
% Since $\psi^2(.) \sim \sics(\rho_t, \frac{1}{\sigma^2})$, we have:
% \begin{equation}
% \Var[\sqrt{\psi(.)}] \approx \frac{\rho_t}{2(\rho_t - 2) (\rho_t - 4) \sigma^2}.
% \label{eqn:variance_appro}
% \end{equation}
% In the next section, we conduct simulation experiments to examine Equation~\ref{eqn:variance_appro} and find it is a reliable approximation. 
% Consequently, we know that $\Var[\sqrt{\psi(.)}]$ decreases at a speed of $O(\frac{1}{\rho_t})$, and the variance in the early stage can be significantly larger than the late stage. 
% For example, as in Figure~\ref{fig:var_approx}, the variance at $\rho = 5$ is over $100$ times larger than the variance at $\rho = 500$.

% \subsection{Variance Rectification}
% In order to assure the adaptive learning rate has consistent variance, we use the following term to rectify the variance of $\rho_t$ as the minimal variance (of $\rho_\infty$): 
% \begin{equation}
% r_t = \sqrt{\frac{\Var[\psi(.)]|_{\rho = \rho_t}}{\Var[\psi(.)]|_{\rho = \rho_\infty}}} = \sqrt{\frac{(\rho_t - 4)(\rho_t - 2)\rho_\infty}{(\rho_\infty - 4)(\rho_\infty - 2)\rho_t}}.
% % \label{eqn:analytic-rectification}
% \end{equation}
% Notice that, since the analytic form is not numerically stable, here we use the first order approximation to derive the rectification term. 
% Applying our rectification term to Adam, our final algorithm is summarized in Algorithm~\ref{algo:cadam}. 
% Specifically, when the length of the approximated SMA is less or equal than 4, the variance of the adaptive learning rate is intractable and the adaptive learning rate is disabled. 
% Otherwise, we will calculate the variance rectification term and update parameters with the adaptive learning rate.

% \begin{algorithm}[t]
% \DontPrintSemicolon
% \KwIn{$\{\alpha_t\}_{t = 1}^T$: step size, 
% $\{\beta_1, \beta_2\}$: decay rate to calculate moving average and moving 2nd moment,
% $\theta_0$: initial parameter, $f_t(\theta)$: stochastic objective function.}
% \KwOut{$\theta_t$: resulting parameters}
% $m_0, v_0 \gets 0, 0$ (Initialize moving 1st and 2nd moment)\;
% $l_\infty \gets 2/(1 - \beta_2) - 1$ (Compute the maximum length of the approximated SMA)\;
% \While{$t = \{1, \cdots, T\}$}{
%     $g_t \gets \Delta_{\theta} f_t(\theta_{t - 1})$ (Calculate gradients w.r.t. stochastic objective at timestep t)\;
%     $v_t \gets \beta_2 v_{t - 1} + (1 - \beta_2) g_t^2$ (Update exponential moving 2nd moment)\;
%     $m_t \gets \beta_1 m_{t - 1} + (1 - \beta_1) g_t$ (Update exponential moving 1st moment)\;
%     $\widehat{m_t} \gets m_t / (1 - \beta_1^t)$ (Compute bias-corrected moving average)\;
%     $l_t \gets l_\infty - 2t\beta_2^t/(1 - \beta_2^t)$(Compute the length of the approximated SMA)\;
%     \uIf{the variance is tractable, \ie, $l_t > 4$}{
%         $\widehat{v_t} \gets \sqrt{v_t / (1 - \beta_2^{t})}$ (Compute bias-corrected moving 2nd moment)\;
%         $r_t \gets \sqrt{\frac{(l_t - 4)(l_t - 2)l_\infty}{(l_\infty - 4)(l_\infty - 2)l_t}}$ (Compute the variance rectification term)\;
%         $\theta_t \gets \theta_{t-1} - \alpha_t r_t \widehat{m_t} / \widehat{v_t}$ (Update parameters with adaptive momentum)\;
%     }
%     \Else{
%         $\theta_t \gets \theta_{t-1} - \alpha_t \widehat{m_t}$ (Update parameters with un-adapted momentum)\;
%     }
% }
% \Return{$\theta_T$}
% \caption{Rectified Adam. All operations are element-wise. }
% \label{algo:cadam}
% \end{algorithm}

% \subsection{Comparison with Warmup}

% We notice that $r_t$ has a similar form with the heuristic linear warmup, which can be viewed as setting the rectification term as $\frac{min(t, T_w)}{T_w}$. 
% It verifies our intuition that warmup works as a variance reduction technique. 
% Comparing these two strategies, our proposed algorithm skips using the adaptive learning rate when its variance is divergent, thus avoiding undesired instability in the first few updates. 
% Besides, our method does not require an additional hyperparameter (\ie, $T_w$) to control the variance reduction and can automatically adapt to different moving average rule. 
